\documentclass[12pt]{extarticle}

% ===== Encodage & langue =====
\usepackage{fontspec}
\usepackage{polyglossia}
\setdefaultlanguage{french}

% ===== Police Verdana (compiler avec XeLaTeX ou LuaLaTeX) =====
\setsansfont{Verdana}
\renewcommand{\familydefault}{\sfdefault}

% ===== Interligne =====
\usepackage{setspace}
\setstretch{1.1}

% ===== Marges réduites =====
\usepackage{geometry}
\geometry{a4paper, top=1.8cm, bottom=1.8cm, left=2cm, right=2cm}

% ===== Mathématiques =====
\usepackage{amsmath}

% ===== Espacement réduit autour des équations =====
\setlength{\abovedisplayskip}{4pt}
\setlength{\belowdisplayskip}{4pt}
\setlength{\abovedisplayshortskip}{2pt}
\setlength{\belowdisplayshortskip}{2pt}

% ===== Espacement réduit des sections =====
\usepackage{titlesec}
\titlespacing*{\section}{0pt}{1.2ex}{0.6ex}

% ===== Titre du document =====
\title{\textbf{Correction du contrôle séquence 1 version 6}}
\author{}
\date{}

\begin{document}
\maketitle

% --------------------------------------------------
\section*{Exercice 1}
\begin{align*}
14 - 2 - 6 + 27 - 5 &= 12 - 6 + 27 - 5 \\
                     &= 6 + 27 - 5 \\
                     &= 33 - 5 \\
                     &= 28
\end{align*}

% --------------------------------------------------
\section*{Exercice 2}

\begin{align*}
\textbf{a)}
4 + 5 \times 6 &= 4 + 30 \\
               &= 34
\end{align*}

\textbf{b)}
\begin{align*}
52 - 7 \times 5 &= 52 - 35 \\
                &= 17
\end{align*}

\textbf{c)}
\begin{align*}
39 + 21 \div 7 &= 39 + 3 \\
               &= 42
\end{align*}

\textbf{d)}
\begin{align*}
4 + 8 \times 5 + 7 - 10 \div 2 &= 4 + 40 + 7 - 5 \\
                              &= 46
\end{align*}

% --------------------------------------------------
\section*{Exercice 3}
\begin{align*}
3 \times 9 + 2 \times 10 &= 27 + 20 \\
                        &= 47
\end{align*}

La dépense est de 47~€.

% --------------------------------------------------
\section*{Exercice 4}

\textbf{a)} $38 - (9 + 7) \div 2$
\begin{align*}
&= 38 - 16 \div 2 \\
&= 38 - 8 \\
&= 30
\end{align*}

\textbf{b)} $(8 + 5) \times 6 - 9 + 7$
\begin{align*}
&= 13 \times 6 - 9 + 7 \\
&= 78 - 9 + 7 \\
&= 76
\end{align*}

\textbf{c)} $(29 - (8 + 9)) \times 6$
\begin{align*}
&= (29 - 17) \times 6 \\
&= 12 \times 6 \\
&= 72
\end{align*}

\textbf{d)} $7 \times (28 - (9 + 13)) - 8$
\begin{align*}
&= 7 \times (28 - 22) - 8 \\
&= 7 \times 6 - 8 \\
&= 42 - 8 \\
&= 34
\end{align*}

% --------------------------------------------------
\section*{Exercice 5}
\begin{align*}
20 - 6 \times 3 &= 20 - 18 \\
                &= 2
\end{align*}

Une gomme coûte 2~€.

% --------------------------------------------------
\section*{Exercice bonus 1}

\textbf{a)} $8 \times 6 + 13$

\textbf{b)} $7 \times (9 + 4)$

% --------------------------------------------------
\section*{Exercice bonus 2}

\textbf{a)} La somme du produit de 8 par 7 et de 6.

\textbf{b)} Le produit de la somme de 4 et de 9 par 7.

\end{document}